% Insert as A.71.  This is an arithmetic group retraction, not a
% construction of a same-pair inverse limit of the A.65 sections.
\begin{statement}{A.71}{Proved}{An exact retraction from pronilpotent to full metabelian arithmetic groups}
Let $M$ be a profinite metabelian group with a continuous action of
a profinite group $G$, put $D=M'$ and $A=M/D$, and write $M^{(\ell)}$
for its natural maximal pro-$\ell$ quotient.  The natural map
\[
 M\rtimes G\longrightarrow
       \left(\prod_\ell M^{(\ell)}\right)\rtimes G
 \tag{A.71.1}
\]
has a continuous group right inverse over $G$ inducing the identity
on $A=\prod_\ell A^{(\ell)}$.
The right inverse depends on choices of coprime complements.

Apply this to $M=\Pi/\Pi''$ and the action defined by the rational
point-$3$ section.  A fixed actual abelian section of the tripod
lifts to the full geometrically metabelian arithmetic quotient
exactly when each of its coefficient-prime components lifts to the
corresponding full geometrically pro-$\ell$ metabelian quotient.
The same equivalence holds for the image on section classes.
This statement constructs the comparison map and retains the
original arithmetic action on every natural pro-$\ell$ quotient.
\end{statement}
\begin{proof}
All derived subgroups in this proof are closed.  The group $D$ is
profinite abelian, so it has the characteristic decomposition
$D=\prod_\ell D_\ell$ into its primary subgroups.
For a fixed $\ell$ put
\[
 \mathcal M_\ell=M/\prod_{p\ne\ell}D_p,\qquad
 A=A_\ell\times A_{\ell'},\qquad
 \mathcal H_\ell=\operatorname{pr}^{-1}(A_{\ell'})
                                      \subset\mathcal M_\ell .
 \tag{A.71.2}
\]
Here $A_{\ell'}=\prod_{p\ne\ell}A_p$.
The kernel of $\mathcal H_\ell\to A_{\ell'}$ is the abelian
pro-$\ell$ group $D_\ell$.

\emph{Coprime complements with infinitely generated kernels.}
There is a closed complement $C_\ell\simeq A_{\ell'}$ to $D_\ell$
in $\mathcal H_\ell$, and any two such complements are conjugate
by $D_\ell$.  The finite-quotient proof is sufficient here.
In each finite quotient the normal kernel is an abelian
$\ell$-group and the quotient has order prime to $\ell$.
Averaging cochains over that quotient makes its positive-degree
cohomology zero: division by the group order is an automorphism of
the kernel.  Vanishing of $H^2$ splits the extension, and vanishing
of $H^1$ makes its complements conjugate by the kernel.
Passing over all finite quotients gives a complement by compactness
of the inverse system of their finite sets of complements.
A finite collection of quotient conditions is realized in one common
finer finite quotient.  For conjugacy, the sets of elements of
$D_\ell$ conjugating the two complements in a specified finite
quotient are nonempty closed sets with this same finite-intersection
property.  Compactness of $D_\ell$ gives a simultaneous conjugator.
This proves both assertions without finite generation of $D_\ell$.

Conjugation on $D_\ell$ factors through $A$ because $D$ is abelian.
Average this action over the pro-$\ell'$ group $A_{\ell'}$:
\[
 e_\ell(d)=\int_{A_{\ell'}}a(d)\,da,\qquad
 D_{\ell,0}=e_\ell D_\ell,\qquad
 D_{\ell,1}=(1-e_\ell)D_\ell.
 \tag{A.71.3}
\]
The integral means the compatible normalized finite averages.
On every finite invariant quotient of $D_\ell$ the acting group
has order prime to $\ell$, so these averages exist and are compatible.
They define a continuous idempotent on the whole compact module.
In particular
\[
 D_\ell=D_{\ell,0}\oplus D_{\ell,1},\qquad
 D_{\ell,0}=D_\ell^{C_\ell},\qquad
 D_{\ell,1}=\overline{[C_\ell,D_\ell]} .
 \tag{A.71.4}
\]
For the last equality, averaging kills every commutator difference,
whereas $d-e_\ell(d)$ is, in each finite quotient, an integral
$\Z_\ell$ linear combination of such differences.
Taking their closed subgroup proves the converse inclusion.

The averaging operator depends only on the action of the
characteristic group $A_{\ell'}$ on $D_\ell$, not on its chosen
complement.  Naturality of normalized Haar averaging therefore
makes both summands in (A.71.4) invariant under the original
$G$-action.  The centralizer of $C_\ell$ inside $D_\ell$ is
$D_{\ell,0}$.  Complement conjugacy consequently identifies the
space of complements with the torsor
$D_\ell/D_{\ell,0}=D_{\ell,1}$.
Conjugation by $\mathcal M_\ell$ factors on $D_\ell$ through
the abelian group $A$, so it commutes with $e_\ell$.
In particular $D_{\ell,1}$ is normal in $\mathcal M_\ell$.

\emph{The geometric normalizer and its natural pro-prime quotient.}
Let $\mathcal N_\ell$ be the normalizer of $C_\ell$ in
$\mathcal M_\ell$.  It centralizes $C_\ell$: conjugation induces
the identity on the abelian quotient $A$, and the map
$C_\ell\to A_{\ell'}$ is injective.
Also $\mathcal N_\ell\to A$ is surjective.
Indeed, for $m\in\mathcal M_\ell$, the groups
$mC_\ell m^{-1}$ and $C_\ell$ are complements in
$\mathcal H_\ell$; conjugacy by $D_\ell$ changes $m$ to an
element of their normalizer without changing its image in $A$.
Thus
\[
 \mathcal N_\ell=C_\ell\times P_\ell,\qquad
 1\longrightarrow D_{\ell,0}\longrightarrow P_\ell
                     \longrightarrow A_\ell\longrightarrow1,
 \tag{A.71.5}
\]
where $P_\ell$ is the kernel of the projection
$\mathcal N_\ell\to A_{\ell'}$.
It is a pro-$\ell$ group.

Furthermore
\[
 \mathcal M_\ell=D_{\ell,1}\rtimes\mathcal N_\ell,\qquad
                      P_\ell\simeq M^{(\ell)}
 \tag{A.71.6}
\]
under the natural quotient map.
For the first assertion, the normalizer maps onto $A$ and contains
$D_{\ell,0}$, while its intersection with $D_{\ell,1}$ is trivial.
For the second, the resulting map
$\mathcal M_\ell\to\mathcal N_\ell/C_\ell=P_\ell$ is a pro-$\ell$
quotient.  Every homomorphism from $\mathcal M_\ell$ to a pro-$\ell$
group kills $C_\ell$ and hence kills
$D_{\ell,1}=\overline{[C_\ell,D_\ell]}$.
It therefore factors through this map.
The other primary subgroups of $D$ were already killed in every
pro-$\ell$ quotient of $M$, so this is exactly its natural maximal
pro-$\ell$ quotient.  In particular the identification in (A.71.6)
does not choose a new action on $M^{(\ell)}$.

\emph{The actual arithmetic correction.}
For $g\in G$ there is a unique $\delta_\ell(g)\in D_{\ell,1}$ with
\[
          g(C_\ell)=\delta_\ell(g)C_\ell\delta_\ell(g)^{-1}.
 \tag{A.71.7}
\]
The map is continuous.  One direct verification is that its graph
is a closed subset of the compact space $G\times D_{\ell,1}$:
the conjugacy condition is closed after every finite quotient.
Its projection to $G$ is bijective by existence and uniqueness,
hence is a homeomorphism.

Since $D_{\ell,1}$ is abelian and $G$-stable, applying (A.71.7)
twice gives the exact cocycle identity
\[
             \delta_\ell(gh)=\delta_\ell(g)+g\delta_\ell(h).
 \tag{A.71.8}
\]
If $s_0(g)=(1,g)$ is the original split arithmetic section, then
\[
                       s'_0(g)=(-\delta_\ell(g),g)
 \tag{A.71.9}
\]
is an actual continuous section of $\mathcal M_\ell\rtimes G$
which normalizes $C_\ell$.
It consequently normalizes $\mathcal N_\ell$ and $P_\ell$.
The action on $C_\ell$ is exactly the action on $A_{\ell'}$.
The action on $P_\ell\simeq M^{(\ell)}$ is also the original
natural one: (A.71.4) shows that $\delta_\ell(g)$ dies in the
natural pro-$\ell$ quotient, so projection sends (A.71.9) to
the original section in $M^{(\ell)}\rtimes G$.

We obtain a continuous arithmetic embedding
\[
 \bigl(A_{\ell'}\times M^{(\ell)}\bigr)\rtimes G
                   \longrightarrow\mathcal M_\ell\rtimes G
 \tag{A.71.10}
\]
by the geometric identification (A.71.5) and the section (A.71.9).
Its map on the common abelianization $A\rtimes G$ is the identity,
and its projection to $M^{(\ell)}\rtimes G$ is the natural projection
on the second factor.

\emph{Assembling all coefficient primes.}
The natural map $M\to\prod_\ell M^{(\ell)}$ is surjective.
Indeed its image in every finite product of finite
$\ell$-group quotients is a subdirect product.
Its order is divisible by the order of each factor, hence by their
product, because these orders are pairwise coprime.
Thus that image is the full product.  Compactness proves the assertion
for the profinite product.
Its abelianization is the natural identification
$A=\prod_\ell A_\ell$.

There is an actual topological group isomorphism
\[
 M\ \xrightarrow{\sim}\
       \prod\nolimits_{A,\ell}\mathcal M_\ell .
 \tag{A.71.11}
\]
Injectivity follows from the primary decomposition of $D$.
For surjectivity take a compatible tuple with common image $a\in A$,
lift $a$ to $m\in M$, and express its difference in each coordinate
as $d_\ell\in D_\ell$.  The element
$(d_\ell)_\ell\in D$ corrects $m$ to the tuple.
The continuous bijection is a homeomorphism by compactness.
The same argument gives the arithmetic fibre product over
$A\rtimes G$.

From the product $\prod_pM^{(p)}$, map to
$A_{\ell'}\times M^{(\ell)}$ by retaining the $\ell$ factor and
taking abelianizations of every other factor.
Compose this equivariant map with (A.71.10).
All the resulting arithmetic maps have identical projection to
$A\rtimes G$, so (A.71.11) assembles them to a continuous
homomorphism
\[
 \Sigma:\left(\prod_\ell M^{(\ell)}\right)\rtimes G
                                      \longrightarrow M\rtimes G .
 \tag{A.71.12}
\]
Projection back to each natural $M^{(\ell)}\rtimes G$ is the identity.
Thus (A.71.12) is the asserted right inverse of (A.71.1), over $G$
and over the given abelianization.

For the tripod, $M^{(\ell)}$ is the maximal pro-$\ell$
metabelian quotient of its geometric group, by the universal
properties of the two quotients.
Given actual coefficient-prime sections with the prescribed
abelian coordinates, their coordinatewise product is continuous
and is a section of the target of (A.71.1).
Composing with $\Sigma$ gives the required full metabelian section.
The converse is projection.
For section classes, lift any abelian conjugator in a coefficient
quotient to align its abelian representative first; then apply the
actual-section construction.  This proves the image assertion.

The retraction changes the chosen arithmetic base section by the
explicit derived-valued cocycles (A.71.9).
Its effect on every abelian coordinate and every natural pro-$\ell$
quotient has been proved above.
A.65 constructs pairs at each fixed depth; it does not provide,
for one fixed pair, the full pro-$\ell$ metabelian sections required
in this final application.
The exact finite certificate
\path{code/metabelian_retraction_certificate.py} checks the complement
cocycle, its sign, and the arithmetic right-inverse homomorphism
for $S_3$ with its inner $S_3$ action and $A_4$ with its $S_4$ action.
The saved output records 612 cocycle identities and 5328 arithmetic
homomorphism identities, including the nontrivial action on
$A_4/V_4$.
\end{proof}

